\begin{pastexam}[label={px:df3-102-ds-02}]{地方三等資訊處理 102 年 第二題}
（申論；20 分）

定義如下的函數 F：\\
如果 x 是偶數，則 F(x) = x/2;\\
否則 F(x) = F(F(3x + 1))
\begin{itemize}[leftmargin=2.2em]
\item[㈠] 請問 F(11) =？（5 分）
\item[㈡] 請證明對於任何正整數 w，我們都可以在有限時間內計算 F(w)。（提示：每個奇數可以寫成(2i + 1)2$^\text{k}$
– 1 的形式，再採用數學歸納法來證明。）（15 分）
\end{itemize}
\end{pastexam}

\begin{pastsoln}
答案：㈠ $F(11) = 13$。 ㈡ 對奇數 $w = (2i+1)2^k - 1$ 的 $k$ 做歸納，證明計算會在有限步內結束，而且
$F(w) = \bigl(3^k(2i + 1) - 1\bigr)/2$。

\textbf{㈠} 依定義一層一層展開（偶數直接除以 $2$）：
\begin{align*}
F(11) &= F(F(34)) = F(17) && (34 \text{ 是偶數}, F(34) = 17)\\
F(17) &= F(F(52)) = F(26) = 13 && (F(52) = 26,\ 26 \text{ 是偶數}).
\end{align*}
所以 $F(11) = 13$。

\textbf{觀察。} $x$ 是奇數時 $3x + 1$ 是偶數，$F(3x + 1) = (3x + 1)/2$，所以
\[
x \text{ 為奇數} \ \Longrightarrow\ F(x) = F\!\left(\frac{3x + 1}{2}\right). \tag{$*$}
\]
也就是說，計算 $F(x)$ 時，每遇到一次奇數，就先花一次偶數的計算（算 $F(3x+1)$），再把問題換成算
$F\bigl((3x+1)/2\bigr)$。 只要證明這個「換成」的動作有限次之後會碰到偶數，計算就會結束。

\textbf{㈡ 證明。} 任何奇數 $w$，$w + 1$ 是偶數，可以唯一寫成 $w + 1 = (2i + 1)2^k$，$i \ge 0$、$k \ge 1$（$2^k$ 是
$w + 1$ 裡 $2$ 的最高次方）。 所以 $w = (2i + 1)2^k - 1$。

\textbf{命題 $P(k)$}：對所有 $i \ge 0$，$F\bigl((2i+1)2^k - 1\bigr)$ 可以在有限步算完，而且值是
$\dfrac{3^k(2i + 1) - 1}{2}$。

\textbf{基底 $k = 1$}：$w = 2(2i + 1) - 1 = 4i + 1$ 是奇數。 由 $(*)$，
\[
F(w) = F\!\left(\frac{3(4i + 1) + 1}{2}\right) = F(6i + 2),
\]
$6i + 2$ 是偶數，$F(6i + 2) = 3i + 1$。 總共只算了 $F(3w + 1)$、$F(6i + 2)$ 兩次偶數的情況，有限步，而且
$3i + 1 = \dfrac{3(2i + 1) - 1}{2}$，$P(1)$ 成立。

\textbf{歸納步驟}：設 $k \ge 2$，且 $P(k - 1)$ 成立。 $w = (2i + 1)2^k - 1$ 是奇數，由 $(*)$，
\[
F(w) = F\!\left(\frac{3w + 1}{2}\right), \qquad
\frac{3w + 1}{2} = \frac{3(2i+1)2^k - 2}{2} = 3(2i + 1)\,2^{k-1} - 1 .
\]
$3(2i + 1)$ 是奇數，記為 $2i' + 1$（$i' = 3i + 1$），所以 $\dfrac{3w+1}{2} = (2i' + 1)2^{k-1} - 1$，正好是
$P(k-1)$ 處理的形式。 由歸納假設，它能在有限步算完，值是
\[
\frac{3^{k-1}(2i' + 1) - 1}{2} = \frac{3^{k-1} \cdot 3(2i+1) - 1}{2} = \frac{3^k(2i+1) - 1}{2}.
\]
$F(w)$ 只比它多算一次 $F(3w + 1)$（偶數，一步），仍是有限步。 $P(k)$ 成立。

由數學歸納法，$P(k)$ 對所有 $k \ge 1$ 成立，也就是每個奇數 $w$ 都能在有限時間內算出 $F(w)$。 偶數 $w$ 則直接
$F(w) = w/2$，一步就完成。 所以對任何正整數 $w$，$F(w)$ 都能在有限時間內算完。 \hfill$\blacksquare$

\textbf{補充。} 從證明可以數出步數：奇數 $w = (2i+1)2^k - 1$ 要做 $k$ 次 $(*)$，每次多算一次偶數，最後再算一次
偶數，共 $2k + 1$ 次呼叫；$k \le \lg(w + 1)$，所以是 $O(\log w)$。 用 ㈠ 驗證：$11 + 1 = 12 = 3 \cdot 2^2$，
$i = 1$、$k = 2$，$F(11) = (3^2 \cdot 3 - 1)/2 = 13$，和直接展開一致。 這題長得像著名的 $3x + 1$ 問題，但那個
問題是反覆做 $x \mapsto 3x + 1$ 或 $x/2$ 直到變成 $1$，至今沒有證明一定會停；這裡的 $F$ 遇到偶數就直接回傳
$x/2$，所以能證明會停。 本題屬於遞迴定義的函數與數學歸納法（資結書附錄~A）。
\end{pastsoln}
